\chapter*{Streszczenie pracy}
\label{cha:streszczenie-pl}

\makeatletter

\addcontentsline{toc}{section}{\textbf{Streszczenie pracy}}

\begin{center}
	\noindent \Large \textbf{\@titlePL}
\end{center}

\vspace*{8mm}
W dobie dynamicznego rozwoju sztucznej inteligencji, systemy klasy asystentów głosowych i tekstowych (takie jak Google Assistant, Amazon Alexa czy Apple Siri) stały się integralną częścią codziennego życia. Mimo swojej popularności, komercyjne rozwiązania tego typu charakteryzują się istotnymi ograniczeniami wynikającymi z zamkniętości ich ekosystemów. Integracja nowych urządzeń lub dodanie niestandardowych funkcjonalności wymaga z reguły korzystania z dedykowanych pakietów SDK, przechodzenia przez rygorystyczne procesy certyfikacji oraz całkowitego uzależnienia się od infrastruktury chmurowej zewnętrznego dostawcy. Skutkuje to brakiem interoperacyjności pomiędzy urządzeniami różnych producentów, a także rodzi obawy dotyczące prywatności. Co więcej, istniejące systemy nie oferują natywnego mechanizmu grupowania różnorodnych urządzeń we wspólną, spójną przestrzeń sterowaną za pomocą pojedynczego polecenia.

Odpowiedzią na wyżej wymienione problemy jest niniejsza praca inżynierska, której głównym celem było zaprojektowanie oraz implementacja rozproszonego, modularnego i w pełni otwartego systemu asystenta AI. Skonstruowane rozwiązanie umożliwia intuicyjne sterowanie heterogenicznym zbiorem urządzeń końcowych (m.in. komputerami osobistymi, smartfonami czy smartwatchami) przy użyciu poleceń wydawanych w języku naturalnym. Fundamentalnym założeniem projektowym była otwartość architektury, pozwalająca na łatwe i dynamiczne rozszerzanie możliwości systemu o obsługę nowych urządzeń oraz dziedzin bez konieczności ingerencji w jego kod źródłowy.

Architektura opracowanego systemu opiera się na modelu klient-serwer. Serwer centralny, zaimplementowany w języku C\#, został zaprojektowany zgodnie z paradygmatem modular monolith z rygorystycznym zastosowaniem metodyki Domain-Driven Design (DDD) oraz wzorca CQRS (Command Query Responsibility Segregation). Taki podział architektoniczny pozwolił na wyraźne wyznaczenie granic kontekstów (bounded contexts) dla poszczególnych modułów systemu (m.in. zarządzanie połączeniami, grupami, wtyczkami i przetwarzaniem zapytań). Ograniczyło to narzut operacyjny charakterystyczny dla rozproszonych architektur opartych na mikroserwisach, przy jednoczesnym zachowaniu wysokiej spójności i podatności na dalszy rozwój. Za dwukierunkową, niskoopóźnieniową komunikację pomiędzy klientami a serwerem odpowiada biblioteka SignalR, zapewniająca niezawodny transport danych w czasie rzeczywistym.

Kluczowym elementem warunkującym rozszerzalność systemu jest innowacyjna architektura pluginów po stronie klienckiej. Funkcjonalności specyficzne dla danego urządzenia są dostarczane w postaci dynamicznie ładowanych bibliotek .dll. W momencie startu, aplikacja kliencka ładuje dostępne moduły do pamięci, po czym rejestruje ich metadane w serwerze centralnym, pozostawiając sam kod wykonawczy wyłącznie na urządzeniu użytkownika. Dzięki temu serwer pełni rolę bezstanowego orkiestratora, nieposiadającego wiedzy o implementacji poszczególnych operacji.

Centralnym bytem logicznym zaprojektowanego rozwiązania są grupy urządzeń. Mechanizm ten pozwala użytkownikowi na łączenie wielu fizycznych maszyn na podstawie wygenerowanego kodu dostępu, tworząc ujednoliconą przestrzeń operacyjną. Gdy użytkownik wydaje polecenie (prompt), serwer agreguje metadane wszystkich aktywnych pluginów we wskazanej grupie i konstruuje rozszerzony kontekst dla dużego modelu językowego (LLM). Integracja z modelami AI została zrealizowana przy użyciu biblioteki Microsoft Semantic Kernel, która stanowi warstwę abstrakcji, uniezależniając logikę biznesową od konkretnego dostawcy LLM. Dzięki zastosowaniu mechanizmu function calling, model poddaje analizie intencję użytkownika i zwraca optymalny zestaw funkcji do wykonania. Po autoryzacji operacji przez użytkownika, serwer dystrybuuje polecenia do docelowych klientów, gdzie są one wykonywane natywnie.

Opracowany system stanowi otwarty ekosystem klasy smart assistant. Zaprojektowana architektura skutecznie rozwiązuje problem interoperacyjności i rozszerzalności, pozwalając na współpracę wielu zróżnicowanych urządzeń. Wykorzystanie nowoczesnych wzorców inżynierii oprogramowania w połączeniu z potęgą dużych modelów językowych umożliwiło stworzenie narzędzia, które stanowi skalowalną alternatywę dla rozwiązań komercyjnych, zapewniając użytkownikowi niezależność i kontrolę nad własnym środowiskiem cyfrowym.
